% table_4.tex
% 偏振特征检测方法量化对比 — 对应 Plan B 柱状图
% 编译: xelatex
% 依赖: \usepackage{booktabs,multirow,amsmath}

\begin{table*}[t]
\centering
\caption{不同特征检测方法在偏振通道上的特征跟踪量化结果}
\label{tab:fp_metrics}
\footnotesize
\setlength{\tabcolsep}{3.5pt}
\begin{tabular}{llccccccccc}
\toprule
& & \multicolumn{3}{c}{\bfseries 明亮光照 (94--205~lux)} & \multicolumn{3}{c}{\bfseries 中等光照 (2.6--18.6~lux)} & \multicolumn{3}{c}{\bfseries 暗光 (0.9--4.8~lux)} \\
\cmidrule(lr){3-5} \cmidrule(lr){6-8} \cmidrule(lr){9-11}
\multirow{-2}{*}{\bfseries 检测方法} & \multirow{-2}{*}{\bfseries 通道} & 内点数$\uparrow$ & 内点比例$\uparrow$ & 时间/ms$\downarrow$ & 内点数$\uparrow$ & 内点比例$\uparrow$ & 时间/ms$\downarrow$ & 内点数$\uparrow$ & 内点比例$\uparrow$ & 时间/ms$\downarrow$ \\
\midrule

% ---- FAST ----
\multirow{3}{*}{FAST}
    & P0 (DoP)     & 73.4  & \textbf{0.943} & \textbf{1.0}  & 1.5   & 0.370 & \textbf{0.8}  & 0.0   & 0.036 & \textbf{0.4} \\
    & P1 (sinAoP)  & 121.9 & \textbf{0.910} & \textbf{1.1}  & 1.0   & 0.299 & \textbf{0.7}  & 0.0   & 0     & \textbf{0.3} \\
    & P2 (cosAoP)  & \textbf{304.6} & \textbf{0.911} & \textbf{1.6}  & 20.1  & \textbf{0.878} & \textbf{1.1}  & 1.9   & 0.465 & \textbf{0.9} \\
\cmidrule{2-11}

% ---- GFTT ----
\multirow{3}{*}{GFTT}
    & P0 (DoP)     & \textbf{152.7} & 0.909 & 3.8  & \textbf{62.8} & 0.710 & 4.2  & \textbf{37.2} & \textbf{0.622} & 5.3 \\
    & P1 (sinAoP)  & \textbf{188.4} & 0.839 & 4.0  & \textbf{24.4} & \textbf{0.653} & 4.7  & \textbf{11.5} & \textbf{0.549} & 6.3 \\
    & P2 (cosAoP)  & 213.2 & 0.882 & 3.9  & \textbf{53.2} & 0.781 & 3.5  & \textbf{20.6} & \textbf{0.650} & 4.6 \\
\cmidrule{2-11}

% ---- SuperPoint ----
\multirow{3}{*}{SuperPoint}
    & P0 (DoP)     & 25.5  & 0.898 & 9.9  & 12.9  & \textbf{0.762} & 10.6 & 5.0   & 0.555 & 10.2 \\
    & P1 (sinAoP)  & 23.0  & 0.847 & 9.9  & 2.4   & 0.361 & 10.6 & 0.1   & 0.032 & 10.2 \\
    & P2 (cosAoP)  & 25.9  & 0.885 & 9.9  & 10.3  & 0.720 & 10.6 & 2.7   & 0.413 & 10.2 \\

\midrule
\end{tabular}

\vspace{4pt}
\begin{minipage}{\textwidth}
\footnotesize
\textbf{注：}
内点数（Avg Inlier）为帧平均 RANSAC 基础矩阵内点数，越高越好，反映后端可用的几何一致特征匹配数量.
内点比例（Inlier Ratio）为内点数与匹配数之比，反映检测方法的匹配精度.
时间为一帧内偏振通道提取-导向滤波预处理-特征点检测-特征点匹配-RANSAC的全流程耗时.
\end{minipage}
\end{table*}
